% !TeX program = xelatex
\documentclass{jos}
\makeatletter
\renewcommand{\jos@topmatter}{\noindent 天津大学NLP平时作业\hfill 第8队个人章节预览\par}
\renewcommand{\jos@citationbox}[2]{}
\makeatother
\fancyhead[RO]{\joskaiti Text2Mem论文阅读\quad ngocnamd93-spec}
\setlength{\headheight}{16pt}
\begin{document}
\Year{2026}
\BeginPage{1}
\CLC{TP18}
\title{Text2Mem：统一记忆操作语言与状态验证}{Text2Mem论文阅读}
\entitle{Text2Mem: Unified Memory Operations and State Verification}{Text2Mem Reading}
\author{ngocnamd93-spec}{ngocnamd93-spec}{}
\AuthorMark{ngocnamd93-spec}
\address{天津大学，第8队}{Tianjin University, Team 8}
\abstract{Text2Mem通过统一的记忆操作规范（Schema）连接自然语言指令与底层存储引擎。本文梳理其三阶段十二种操作、五字段Schema以及校验、解析与适配执行流水线，并结合结构匹配度、执行成功率与状态断言满足率分析主流模型的表现。实验表明，现有模型虽能生成语法合规的可执行操作，但在复杂存储操作与多步工作流的状态准确性上仍存明显瓶颈，凸显了后置状态核验的必要性。}
\keywords{记忆操作；结构化规范；状态验证；论文阅读}
\maketitle
\section{Text2Mem：记忆操作的统一表达与状态验证}
\label{sec:text2mem}

Text2Mem\upcite{wang-etal-2026-text2mem}面向自然语言驱动的记忆管理，提出统一的记忆操作规范（Schema）与配套评测基准。系统将更新、合并、降级或定期清理等自然语言指令解析为结构化操作对象，映射至底层存储执行，并通过后置状态断言核验执行效果。

\subsection{研究问题与适用范围}

长期交互不仅需要检索历史信息，还需对记忆的优先级、有效期、合并与删除进行精细控制。日常指令往往存在操作语义歧义，例如“暂时不要再提午餐”通常意指下调显著性权重，而非永久删除记录。若缺乏明确的操作类型、作用范围与时间参数约束，系统即便成功调用接口，也可能误改无关记忆。Text2Mem因此在自然语言指令与底层存储引擎之间引入结构化操作中间层，将记忆生命周期语义与具体后端实现解耦（论文主实验基于SQL/SQLite环境评测）。

\subsection{操作集合、Schema与执行流程}

论文围绕编码（ENC）、存储（STO）和检索（RET）三个阶段定义了十二种操作：编码阶段含Encode；存储阶段含Update、Label、Promote、Demote、Merge、Split、Delete、Lock和Expire九种操作；检索阶段含Retrieve与Summarize。其中，Encode组织内容、语义标签、来源与治理属性（支持延迟生成向量）；Promote与Demote调整记忆显著性（Demote支持归档降权而不删除）；Delete处理软／硬删除，Expire设置到期规则，三者在生命周期语义上严格区分。

操作规范包含\texttt{stage}、\texttt{op}、\texttt{target}、\texttt{args}和\texttt{meta}五个字段，依次对应阶段、动作、目标、专属参数与执行元数据。不同操作对应特定的必填与跨字段约束：例如Encode新建记录时无需先选定已有目标；针对存量记忆的操作则需在\texttt{target}中从\texttt{ids}、\texttt{filter}、\texttt{search}、\texttt{all}四类选择器中择一明确范围，并受锁定保护与批量写入权限限制。

执行流水线中，Validator首先校验结构完整性与跨字段约束；Parser将合法实例解析为类型化操作并归一化时间与优先级参数；Adapter最终将其映射为后端数据库动作，返回记录受影响条目和状态变化的\texttt{ExecutionResult}。多步请求可表示为保留依赖顺序的\texttt{schema\_list}依次执行。

\begin{samepage}
例如，针对“提高本周课设进度笔记的重要性，其他记录不变”，系统对条目17执行STO阶段的Promote操作，将\texttt{ids=[17]}的\texttt{weight}设为0.9，并在执行后同时核验目标权重、原有内容与非目标条目，确保仅改动指定范围。

\end{samepage}

\begin{samepage}
\subsection{基准、指标与主要结果}

评测基准由自然语言指令、前置操作建立的初态、参考Schema和后置SQL状态断言组成，覆盖中英双语、直接与间接表达、单步操作与多步工作流，并在统一提示模板下评测了GPT-4o、Qwen2.5-72B-Instruct、Claude-4-Sonnet、Gemini-2.5-Pro与GPT-5五种模型。

\end{samepage}

评测包含三项指标：结构匹配度（SMA）用编辑距离或余弦相似度衡量生成Schema与参考Schema的一致性；执行成功率（ESR）统计通过校验且无运行错误的实例比例；执行匹配率（EMR）统计执行后满足的SQL状态断言占全部期望断言的比例：
\begin{equation}
    \mathrm{EMR}=\frac{\sum_i\sum_{a\in A_i}\mathbb{I}(a)}{\sum_i|A_i|}.
\end{equation}
其中$A_i$为第$i$个实例的断言集合，断言$a$满足时$\mathbb{I}(a)=1$，否则为0。由于EMR按断言汇总，包含更多检查项的复杂实例对总体得分权重更高。

\begin{table}[!htbp]
    \centering
    \caption{Text2Mem总体评测指标}
    \label{tab:text2mem-main}
    \small
    \begin{tabular}{lrrrr}
        \toprule
        模型 & SMA$_{\rm lev}$ & SMA$_{\rm cos}$ & ESR & EMR \\
        \midrule
        GPT-4o & 0.6776 & 0.8690 & 0.9484 & 0.2532 \\
        Qwen2.5-72B-Instruct & 0.6195 & 0.8349 & 0.9510 & 0.2535 \\
        Claude-4-Sonnet & 0.6230 & 0.8475 & 0.9594 & 0.3132 \\
        Gemini-2.5-Pro & 0.6491 & 0.8560 & 0.9479 & 0.2897 \\
        GPT-5 & 0.6115 & 0.8396 & 0.9662 & 0.2826 \\
        \bottomrule
    \end{tabular}
\end{table}

表\ref{tab:text2mem-main}显示，各模型的ESR均达94.79\%--96.62\%（GPT-5最高），但EMR仅为25.32\%--31.32\%（Claude-4-Sonnet最高）。这一落差说明当前模型虽能稳定生成语法合规、可执行的操作指令，但在准确锁定目标范围并达成预期存储状态上仍有明显瓶颈。

在复杂操作与多步工作流上差距尤为突出：五个模型的Split EMR均为0，Merge仅GPT-4o取得0.0465，其余为0；GPT-5的单步操作EMR为0.2930，而在多步工作流中尽管ESR为1.0000，EMR却降至0.1190，表明跨步状态依赖传递是主要难点。

\subsection{方法局限与课设启发}

Text2Mem将结构校验（ESR）与后置状态核验（EMR）分离，清晰揭示了“指令可执行”与“状态改得对”之间的差距。不过，现有实验主要在离线SQL环境中评测，尚未接入长程对话循环检验状态更新对后续问答的长期影响；若状态断言未充分覆盖非目标记录，也可能漏检误改旁支数据的副作用。在课程设计中，为StructMem补充计划变更与失效维护时，可借鉴Text2Mem将模糊变更要求显式化为Update、Demote或Expire操作，并在执行后同时核验目标事件与无关历史条目。

% 原文定位：原文第1--3页，图1；第5、7页，\S3.3、\S4.3；第4页表1，第11页附录A；第4--5页，\S3.2；第3页图2、第5--7页\S3.3--4.1及图3；第12页附录B.1.1；后两项为本文的检查建议；第6--8页，\S4；第7--8页，\S4.1、\S4.3；第8、15页；第9页Limitations。

\bibliographystyle{jos}
\bibliography{references}
\end{document}
